\chapter{Cardiac Ablation}

\section*{Overview}
Cardiac ablation uses heat (radiofrequency) or cold (cryoablation) energy to scar tiny areas in the heart that are creating or conducting irregular electrical signals, thereby restoring normal sinus rhythm.

\section*{Alternative Names}
Catheter ablation, Radiofrequency ablation, Cryoablation

\section*{Principle}
Catheters with electrode tips are advanced into the heart to map abnormal electrical pathways, then deliver radiofrequency energy or extreme cold to create small scars. The scars interrupt the reentrant circuits that cause arrhythmias, restoring normal rhythm.
\section*{History}
Pioneered in the 1980s using direct current shocks. The technology evolved to radiofrequency catheter ablation in the 1990s, becoming the primary curative treatment for many arrhythmias.

\section*{Why It Is Done}
Ordered to treat symptomatic arrhythmias like atrial fibrillation, atrial flutter, or supraventricular tachycardia (SVT) when medications are ineffective or not tolerated.

\section*{Is This a Primary or Secondary Test or Procedure}
Primary curative treatment for specific tachyarrhythmias (SVT, WPW, atrial flutter).

\section*{How It Is Performed}
Performed under sedation or general anesthesia. Diagnostic catheters are inserted through blood vessels in the groin and threaded to the heart. Mapping identifies abnormal pathways, and the ablation catheter delivers energy to isolate or destroy the targeted tissue.

\section*{Preparation}
Fast for 6-8 hours before the procedure. Stop antiarrhythmic drugs 3-5 half-lives before the procedure as directed by your electrophysiologist.

\section*{Contraindications and Precautions}
Active infection, intracardiac thrombus (must be ruled out via transesophageal echo), or severe uncorrected bleeding disorders.

\section*{Risks}
Cardiac tamponade, groin hematoma, femoral pseudoaneurysm, stroke, transient ischemic attack (TIA), heart block requiring a permanent pacemaker, or phrenic nerve injury.

\section*{Aftercare and Recovery}
Lie flat for 4-6 hours post-procedure to prevent groin bleeding. Monitor the puncture site. Avoid strenuous activity and heavy lifting for 1 week.

\section*{Outcome and Follow-up}
Success is defined by the inability to induce the target arrhythmia in the electrophysiology lab post-ablation.

\section*{Expected Results}
Termination of the abnormal pathway; restoration and maintenance of normal sinus rhythm.

\section*{Follow-up}
ECG monitoring, Holter monitoring, and cardiology check-up in 1 month.

\section*{When to Seek Medical Care}
Follow the procedure team's specific instructions. Seek urgent medical assessment for severe or worsening pain, uncontrolled bleeding, fever or signs of infection, trouble breathing, chest pain, fainting, new weakness or confusion, inability to urinate, or any rapidly worsening symptom. The appropriate warning signs depend on the procedure and the patient's medical condition.

\section*{References}
\begin{enumerate}
  \item MedlinePlus. Cardiac Ablation. U.S. National Library of Medicine; 2025.
  \item Current Medical Diagnosis and Treatment. McGraw-Hill; 2025.
\end{enumerate}

\clearpage
